\item A study was conducted among individuals with chronic pain to determine how effective acupuncture is in relieving pain. A random sample of 15 subjects were chosen and their sensory rates were measured. Assume that the measurements on sensory rates follow a  normal distribution. The data from the study are given below:


\begin{verbatim}
  8.6  9.4  7.9  6.8  8.3  7.3  9.2  9.6  8.7  11.4  10.3  5.4  8.1  5.5  6.9
\end{verbatim}

\begin{enumerate}
      \item What is the value of $\bar{x}$, the sample mean (Use SAS, R, JMP or MS Excel to calculate)?  (Round your answer to the nearest 2 decimal places) \\ %\red{$\bar{x} = 25.44$} 
     \item What is the value of $s$, the sample standard deviation (Use SAS, R, JMP or MS Excel to calculate)? \ (Round your answer to the nearest 2 decimal places) \\ %\red{$s = 13.677$} 
     \item Report the degrees of freedom associated with the t-distribution.\\
     \item What is the value of $t_{\frac{\alpha}{2},n-1}$ for the 95\% confidence interval for the mean sensory rate? (The value obtained from the t table, accurate to the nearest 3 decimal places) \\ %\red{$t^* = 2.262$} 
      \item What is the margin of error for the 95\% confidence interval for the mean sensory rate? \ (Round your answer to the nearest 2 decimal places) \\ %\red{$m=2.262 \times \frac{13.677}{\sqrt{10}}=9.78$} 
      \item Calculate the 95\% confidence interval for the mean sensory rate.
      \begin{enumerate}
      \item What is the value of the lower bound of the 95\% confidence interval? \ (Round your answer to the nearest 2 decimal places) %\red{$15.66$} 
      \item What is the value of the upper bound of the 95\% confidence interval?  \ (Round your answer to the nearest 2 decimal places)\\ %\red{$35.22$} 
%\end{enumerate}
\end{enumerate}
\item  Interpret the confidence interval you constructed above in the context of the problem.\\

%\red{We are 95\% confident that the true average cost of an order during lunchtime on Wednesdays is between \$16.76 and \$34.12.}
      \item Assume that the standard deviation of all the sensory rates is known to be 1.67. %that is $\sigma = 1.67$. % ,i.e, \textbf {$\sigma$ is known}.
       Calculate the 95\% confidence interval for the mean sensory rate under this situation. 
      \begin{enumerate}
      \item What is the value of the lower bound of the 95\% confidence interval? (Round your answer to the nearest 2 decimal places) 
      \item What is the value of the upper bound of the 95\% confidence interval? (Round your answer to the nearest 2 decimal places)\\
      \end{enumerate}
      \item Is the 95\% confidence interval constructed in part(f) \textbf{wider} or \textbf{narrower} than the 95\% confidence interval in part(h)? Explain why.\\ %\red{$35.22$} 

  %\red{Our confidence interval is $25.44 \pm 2.262 \frac{13.677}{\sqrt{10}}$ which gives $(15.66, 35.22)$. Thus, we are 95\% confident that the true average cost of an order during lunchtime on Wednesdays is between \$15.66 and \$35.22.}\\
%\end{enumerate}


\end{enumerate}

